% mage-auto-heading: what-must-transfer
% mage-auto-heading: model-based-engineering-and-digital-twins
% mage-auto-heading: professional-knowledge-and-scientific-work
% mage-auto-heading: where-the-generalization-fails
% mage-auto-heading: a-research-prediction
% Authoritative LaTeX manuscript; Markdown is a generated compatibility projection.

% <!-- part-title: The Profession -->

% <!-- chapter-title: Agentic Engineering Beyond Software -->

\chapter{Agentic Engineering Beyond Software}

% <!-- point: mage-composes-two-inheritances-around-delegated-realization | MAGE combines engineering through representations, evidence, and control with AI techniques that extend reasoners through external representations, memory, and tools. -->
MAGE combines two established lines of work around delegated realization. Engineering disciplines developed ways to represent systems, justify consequential claims, and control realization. AI developed ways to extend reasoners through external representations, memory, and tools. \autoref{established-traditions-new-composition} shows how these contributions meet in an agentic production system. Neither tradition is specific to software production.

\begin{figure}[htbp]
\centering
\includegraphics[width=\linewidth]{book/latex/assets/established-traditions-new-composition-bd6313723f.pdf}
\caption{\emph{Established traditions, new composition.} MAGE draws on two long-running traditions. Engineering disciplines developed explicit models, constrained representations, verification, and external control; AI repeatedly extended reasoners through knowledge representation, planning state, memory, structured reasoning, tools, and surrounding machinery. MAGE brings these inheritances together around autonomous engineering work: Modeling makes engineering knowledge explicit, Alignment enforces selected obligations independently, and both meet in the Governed Engineering Environment. Governance conversion preserves recurring engineering judgment in reusable structure within that environment.}
\label{established-traditions-new-composition}
\end{figure}

% <!-- point: what-mage-claims-is-new | MAGE's contribution concerns engineering the surrounding environment when autonomous implementation becomes a primary productive substrate; the broader proposal extends that question to consequential work in other domains. -->
Much recent agentic software engineering asks how increasingly capable agents can perform work within the software-engineering environment they inherit. MAGE asks how that environment itself should be redesigned when autonomous implementation becomes a primary productive substrate. Its contribution lies in making purposeful models, Alignment machinery, and the governed engineering environment objects of engineering, with recurring judgments preserved in structures that subsequent work inherits. The broader question is whether the same approach can govern delegated design, analysis, and decisions outside software.

Software offers unusually favorable conditions for this approach. It is a symbolic artifact whose consequential properties can often be represented precisely. Its behavior remains revisable after construction, frequently at low realization cost. Version control, tests, static analysis, permissions, admission gates, deployment controls, and telemetry provide mature machinery for controlling those changes. Software engineering is the engineering discipline of \textbf{controlled change}~\citep{davis2026sehandbook}: the artifact must remain acceptable as both it and its environment evolve. Representability, revisability, and this controlled-change machinery help explain why software is a promising setting for MAGE.

The evidence in this book comes primarily from software; transfer beyond it remains a hypothesis. Which of software's favorable conditions survive when the result is a physical artifact, a financial commitment, a legal argument, or a scientific claim? The comparison must concern the properties and actions that can be governed, not simply the presence of capable agents in another profession.

\section{What Must Transfer}\label{heading:what-must-transfer}

The first requirement is representability: a purposeful reduction must preserve the distinctions relevant to a consequential decision. Beyond software, the represented object might be a proposed design, the state of a business, the alternatives in a decision, or the evidence supporting a claim. A structural model can preserve the quantities needed to reason about a bridge; a decision model can preserve alternatives, objectives, assumptions, uncertainties, and predicted consequences needed to reason about a business choice. The form must make the domain's question easier to answer.

Representation alone is insufficient. MAGE also depends on correspondence: some credible account of why the representation describes what it claims to describe. Drawings and physical models can be reconciled with measurements, inspections, and the realized artifact. In knowledge work, a representation may claim to describe the relevant state of the world, the evidence available when a decision was made, or the consequences expected to follow from it. Those claims can be checked against source data, independent evidence, and eventually observed outcomes. A model that has drifted from its object can make reasoning easier while making its conclusions worse.

A third condition is governability. At least some consequential actions, outputs, or transitions must be susceptible to checks or controls outside the producing agent. Work may be constrained, routed, admitted, rejected, escalated for human judgment, or subjected to additional evidence requirements. An engineering design might require certification; a financial commitment might require human approval above a threshold. The important distinction is between asking the producing agent to remember an obligation and engineering an environment in which the obligation has an independent path to consequence.

A fourth condition is economy. Models, controls, and correspondence mechanisms cost something to build and maintain. Their value lies in avoiding costs elsewhere: repeated reconstruction, repeated judgment, review, error, delay, or risk. The balance will differ by domain and by task frequency. A representation worth maintaining for work performed ten thousand times may not be worth constructing for a singular judgment.

% <!-- point: unit-of-transfer-is-the-surface | Transfer depends on the properties, evidence, actions, and economics of an engineering surface rather than on a profession as a whole. -->
The conditions need not hold equally across a domain. Strong Modeling may be useful even where few consequential judgments admit independent enforcement. MAGE therefore transfers by engineering surface rather than by profession: the relevant question is which properties can be represented, which claims can be supported by evidence, which actions can be governed, and at what cost. The forms that serve those functions are domain-specific.

\section{Model-Based Engineering and Digital Twins}\label{heading:model-based-engineering-and-digital-twins}

Controlled change is a useful test of what transfers to physical engineering. Software can often be revised repeatedly while tests, permissions, and deployment controls preserve selected obligations. Model-Based Systems Engineering already uses models to support requirements, design, analysis, verification, and validation across a system's lifecycle~\citep{incose2023handbook}. The opportunity for agents is therefore broader than producing an initial design: they could help engineers evaluate and govern changes to a system already in service. Whether that work can be delegated depends on how far the representations and controls reach into the actual change.

The possibility raised in \autoref{7.2-software-rejoins-engineering:heading:model-based-engineering-at-smaller-scales} also matters here: agents may change the scale at which such engineering environments become economical. Smaller organizations could then construct project-specific models and analyses without first developing a comprehensive modeling infrastructure.

Digital twins can support that work by maintaining computational representations in correspondence with physical systems. Observations of the physical system inform the representation, which supports analysis of the system across its lifecycle~\citep{grieves2017digitaltwin}. An agent could use such a twin to investigate observed behavior, evaluate a proposed operating change, or identify the measurements needed to distinguish competing explanations. The value comes from the connection between the represented state and independent observations of the system.

A maintained twin does not by itself supply the rest of a governed change process. Engineers need a representation of the proposed change, analyses capable of evaluating the consequential properties it might affect, and a way to realize the approved change. They also need evidence that the resulting system corresponds to what was approved and controls over who or what may authorize the intervention. A twin may describe current behavior accurately while omitting a failure mode introduced by a proposed modification. Its predictions may be useful without being sufficient evidence for authorizing that modification. Gaps in these capabilities limit the work that agents can responsibly carry through the twin.

Where those gaps can be closed, agents could help adapt physical systems through maintained representations and independent evidence. A digital twin can make a revision easier to reason about and govern; it does not make concrete, steel, machinery, or physical deployment cheap to change. Slow or irreversible realization raises the value of analysis before action and changes the economics of experimentation. Conventional engineering shows the transfer most clearly when the outcome is an artifact. The harder test is whether the same control structure survives when the outcome is a decision or action instead.

\section{Professional Knowledge and Scientific Work}\label{heading:professional-knowledge-and-scientific-work}

In professional knowledge work, acceptance may authorize a consequential decision rather than release an artifact to a separate user. An accounting entry changes a financial record, a legal filing advances a claim, and a scientific analysis supports a conclusion that others may rely on. The transfer criteria ask what can be represented, how it connects to evidence, and where an independent mechanism can govern the action. Accounting, law, and science expose different answers.

% <!-- point: accounting-environment-already-engineered | Double-entry makes balancing checkable; reconciliation, source records, and allocated authority support correspondence and govern selected actions while classification, estimates, and materiality require further judgment. -->
Accounting supplies relatively structured representations and controls. Double-entry bookkeeping makes one obligation explicit: the debits and credits recorded for a transaction must balance. Arithmetic can expose a violation, but balanced entries can still omit a transaction or classify it incorrectly. Ledgers, charts of accounts, reconciliations, and audit trails support correspondence with economic activity; approval limits, segregation of duties, and closing controls govern selected actions. An agent can work through these representations while checks operate independently of its reasoning. Materiality, ambiguous classification, and the reasonableness of an estimate still require evidence and professional judgment beyond whether the books balance.

% <!-- point: law-representation-without-determinization | Legal representations organize facts and authorities without resolving contested interpretation; procedural controls can govern filings while substantive conclusions remain matters of justified judgment. -->
Legal analysis has structured evidence but more contested interpretation. Parties, claims, facts, authorities, procedural requirements, and relationships among precedents can be represented explicitly. Agents can trace propositions to authority and identify missing evidence or conflicting authorities. Deadlines can be monitored, and submissions missing required components can be blocked. Two lawyers working from the same represented facts and authorities may nevertheless responsibly reach different conclusions. Modeling makes their disagreement easier to examine; procedural conformance does not settle credibility, strategy, or the weight of competing interpretations.

% <!-- point: science-tests-the-theory | Scientific representations and procedural checks make inquiry examinable without establishing explanatory merit; reproducibility does not by itself establish discovery, and inquiry may revise its measurement models and assumptions. -->
Scientific work adds the possibility that inquiry must revise the models through which its evidence is interpreted. Hypotheses, experimental designs, datasets, provenance, measurement models, and analysis plans can be represented and inherited. An environment can record preregistration, require departures from a plan to be disclosed, check data integrity, and refuse admission to a computational result until its steps reproduce. But a reproducible calculation can rest on a poor measurement model, and a surprising observation may justify changing the analysis or designing another experiment. These controls establish conditions under which a result was produced. They do not establish that an explanation is compelling or that a discovery has been made.

The comparison separates procedural governability from substantive acceptance. Accounting often supplies precise invariants, but completeness and classification still depend on correspondence with economic activity. Law can preserve facts and authorities while leaving their interpretation contested. Science can govern provenance and reproduction while permitting its representations to be challenged. In each case, recurring checks may justify durable machinery; infrequent or unsettled judgments may call for representations and review whose costs are justified by the particular decision. The relevant balance is between the assurance the environment can establish and the judgment needed to accept what follows from it.

\section{Where the Generalization Fails}\label{heading:where-the-generalization-fails}

% <!-- point: failure-traits-named | Transfer is limited by tacit or embodied knowledge, reflexive effects, context-dependent obligations, missing intervention boundaries or evidence, and poor economics of maintaining representations. -->
The transfer conditions also identify where MAGE may be difficult or inappropriate. Some knowledge is tacit or embodied: the surgeon's feel for tissue, the negotiator's read of a room, the mechanic's ear for a failing bearing. If a representation omits what trained perception contributes, a check over it cannot establish the whole judgment. Other obligations around the work may still be governable.

Some work changes its object by representing it. A disclosed negotiating position, a published trading model, or a community under observation may alter the behavior being described. Correspondence must account for that response. A representation that assumes the object remains unchanged can mislead even if it accurately described the earlier situation.

Some consequential obligations depend on context or disputed evidence. A clinical decision or a use-of-force judgment may turn on distinctions a proposed rule does not preserve. Mechanizing that rule would enforce a proxy for the obligation. Similarly, applying an agreed decision rule does not resolve disagreement about what the evidence warrants. The engineer must identify those gaps before treating compliance as sufficient.

Independent Alignment also needs somewhere to act. If consequence follows immediately from the producing reasoner's action, with no constraint, approval boundary, or opportunity to obtain independent evidence, the environment has little means to intervene. Highly novel work presents an economic limit instead: existing models may be inadequate and recurrence uncertain, so constructing and maintaining durable structure may not repay its cost.

These limits can leave useful room for representation even when independent enforcement remains weak. Modeling can help even where strong Alignment cannot.

\section{A Research Prediction}\label{heading:a-research-prediction}

The generalization makes two testable claims. The first concerns performance: holding agent capability roughly fixed, consequential work should become cheaper, more reliable, or require less human reconstruction when the environment supplies trustworthy task-relevant representations and independently governs obligations that would otherwise be repeatedly inferred. The effect should grow with task scale and recurrence, and shrink as representation, correspondence, or enforcement becomes expensive or unreliable.

The second, more ambitious claim concerns convergence. Independently developed systems for governing consequential agentic work should acquire similar structures despite producing different outcomes: purposeful representations connected to their subjects, independently produced evidence, allocated authority, and mechanisms that preserve recurring judgment for later work. Comparative studies should examine whether these structures emerge, as well as whether introducing them improves performance. Similar architecture alone would not establish economic benefit, and success in one domain would not establish general transfer.

The failure conditions must bound both comparisons. The broader claim should fail if software's unusually symbolic realization substrate turns out to be essential rather than merely favorable. These hypotheses have been developed here from software engineering; testing them elsewhere is future work.
